\newpage

\section{Opis Metodyki SCRUM}

Scrum to zwinna metodyka zarządzania projektami, najczęściej stosowana w tworzeniu oprogramowania, ale coraz częściej także w innych branżach. Jej głównym celem jest dostarczanie wartościowego produktu w krótkich, regularnych cyklach oraz szybkie reagowanie na zmiany wymagań.

Podstawą Scruma są iteracje, nazywane Sprintami, które zazwyczaj trwają od 1 do 4 tygodni. Każdy Sprint kończy się dostarczeniem działającego fragmentu produktu. Praca zespołu opiera się na jasno określonych rolach: Product Owner odpowiada definiuje wizję produktu i jego funkcjonalność, Scrum Master kieruje zespół ku zaimplementowaniu wymagań, a deweloperzy faktycznie te wymagania realizują.

W Scrumie duże znaczenie mają także wydarzenia, takie jak planowanie Sprintu, codzienne krótkie spotkania (Daily Scrum), przegląd Sprintu oraz retrospektywa, podczas której zespół analizuje swoją pracę i szuka sposobów na jej usprawnienie. Wszystkie zadania są gromadzone w Backlogu Produktu,  który jest regularnie aktualizowany i porządkowany według priorytetów zespołu.
